\subsection{Axioms, coefficients, and the five comparison classes}

\begin{claim}{A.1}{Proved}{Definitions and the axioms actually required}
Let $\D:\Cat_{\U}^{\mathrm{coop}}\to\mathrm{CAT}_{\V}$ be a
prederivator. For a small category $A$, let $\LC_\D(A)$ be the full
subcategory of $\D(A)$ on objects whose underlying presheaf
$A^{\op}\to\D(e)$ sends every arrow to an isomorphism. If
$f:A\to B$, restriction gives a functor
$f^*:\LC_\D(B)\to\LC_\D(A)$. We define $\Wlc$ by this functor
being an equivalence and $\Wone$ by its being fully faithful.
Writing $p:A\to e$, $q:B\to e$, we define $\Wc$ by bijectivity of
\[
 \Hom_{\D(B)}(q^*M,q^*N)\longrightarrow
 \Hom_{\D(A)}(p^*M,p^*N)\qquad(M,N\in\D(e)).
 \tag{A.1.1}
\]
If all restrictions have adjoints $f_!\dashv f^*\dashv f_*$, we
define $\Wcoh$, respectively $\Whom$, by invertibility, for every
$F\in\LC_\D(B)$, of the specified comparison
\[
 q_*F\longrightarrow p_*f^*F,
 \qquad
 p_!f^*F\longrightarrow q_!F
 \quad\text{in }\D(e).
 \tag{A.1.2}
\]
The first map is induced by the unit $F\to f_*f^*F$; the second by
the counit $f_!f^*F\to F$. The definitions require no enhancement.
\end{claim}
\begin{proof}
Evaluation and the underlying presheaf commute with restriction, so
restriction preserves the stated full subcategories. Pseudofunctorial
coherence identifies $qf$ with $p$ and supplies (A.1.1).
The adjunctions and their composition isomorphisms supply exactly
(A.1.2). These constructions do not assume that the underlying diagram
functor is full or faithful.
\end{proof}

Unless explicitly stated otherwise, a derivator satisfies the usual four
axioms in presheaf variance: disjoint unions of shapes become products,
including $\D(\varnothing)\simeq e$; point evaluations jointly detect
isomorphisms; every restriction has both adjoints; and both Kan extensions
are computed by comma categories. In particular $b^*u_!$ is computed on
$b\backslash A$ and $b^*u_*$ on $A/b$, with presheaf variance on those
shapes. We do not use the strongness axiom that coherent arrows present
all incoherent arrows, unless a statement says so. These conventions agree
with the definitions used by \citet[\S3.2]{Cisinski2009LocallyConstant}; they are not
Grothendieck's original numerical labels Der~(1)--(6).

\begin{claim}{A.2}{Proved}{The inclusion chain and its two separate last arrows}
For a prederivator with the adjoints of A.1,
\[
 \Wlc\subseteq\Wone\subseteq
 \Wcoh\cap\Whom\subseteq\Wc.
 \tag{A.2.1}
\]
In fact each of $\Wcoh$ and $\Whom$ is separately contained in
$\Wc$. For a derivator,
$\Winfty\subseteq\Wlc$.
\end{claim}
\begin{proof}
The first inclusion is the definition of equivalence of categories.
If $f^*$ is fully faithful on locally constant objects, then, for
$M\in\D(e)$ and $F\in\LC_\D(B)$, it gives bijections on
$\Hom(q^*M,F)$ and on $\Hom(F,q^*M)$. The adjunctions identify these
with the maps obtained by applying $\Hom(M,-)$, respectively
$\Hom(-,M)$, to (A.1.2). The ordinary Yoneda lemma in $\D(e)$ makes
both comparison maps isomorphisms. Taking $F=q^*N$ proves the last
arrows by the same adjunctions. The final assertion is exactly
\citet[Corollary~3.8]{Cisinski2009LocallyConstant}, with $A^{\op},B^{\op}$ if
necessary to match variance. Its proof uses Theorem~3.6 and the
model-category result Theorem~1.3; it applies to derivators, not merely
to derivators known in advance to come from a model category.
\end{proof}

\begin{claim}{A.3}{Proved}{What happens for a map to the point}
For every prederivator and every $A$,
\[
 (A\to e)\in\Wc\quad\Longleftrightarrow\quad
 (A\to e)\in\Wone.
 \tag{A.3.1}
\]
Consequently the claim in the CATHOT addendum that an acyclic-group
projection $BA\to e$ is a constant-coefficient equivalence but is not
1-complete is false.
\end{claim}
\begin{proof}
Every object of $\D(e)=\LC_\D(e)$ is a coefficient on the target.
Full faithfulness of $p^*:\D(e)\to\LC_\D(A)$ is precisely
(A.1.1) with $B=e$, since $\LC_\D(A)$ is full in $\D(A)$.
This argument includes $A=\varnothing$ and does not use acyclicity.
The correctly oriented strictness examples occur in Part~C.
\end{proof}

\subsection{Formal closure and coherence}

\begin{claim}{A.4}{Proved}{Closure properties that follow without descent}
For any prederivator, $\Wc$ and $\Wlc$ contain identities and are
closed under composition, two-out-of-three, and retracts of arrows.
The class $\Wone$ contains identities, is closed under composition
and retracts, and satisfies the cancellation rule
\[
 f\in\Wone,\quad gf\in\Wone\quad\Longrightarrow\quad g\in\Wone.
 \tag{A.4.1}
\]
With adjoints, $\Wcoh$ and $\Whom$ have these latter closure
properties too. The missing two-out-of-three implication for these
three intermediate classes is false even for a represented derivator.
\end{claim}
\begin{proof}
For $\Wc$, (A.1.1) is a natural map of sets for each fixed $M,N$;
bijections have the asserted closure properties. The assignment
$A\mapsto\LC_\D(A)$ is contravariantly pseudofunctorial, and
equivalences of categories have the same properties. A retract of a
functor which is an equivalence is fully faithful because its maps on
Hom sets are retracts of bijections, and essentially surjective by
applying the retraction to an isomorphism exhibiting essential
surjectivity. The same Hom-set argument handles fully faithful functors.
If $f^*$ and $f^*g^*$ are fully faithful, cancellation of the bijective
Hom maps for $f^*$ proves full faithfulness of $g^*$.

For cohomology, the comparison for $gf$ is the composite of the
comparison for $g$ and that for $f$ applied to $g^*F$; the latter
coefficient is locally constant. This proves composition and (A.4.1).
For a retract diagram in the arrow category, pull a target coefficient
back along the target retraction. Its comparison restricts to the
original one, which is therefore a retract of an isomorphism. The
homological proof is identical with the arrows reversed. A.5 supplies
the stated failure.
\end{proof}

\begin{claim}{A.5}{Refuted}{Two-out-of-three for 1-completeness}
The assertion that $\Wone$ is a fundamental localizer for every
derivator is false. The same obstruction applies to $\Wcoh$ and
$\Whom$.
\end{claim}
\begin{proof}
Take the represented derivator of sets. Let $G=C_2=\{1,s\}$,
$s^2=1$, and let $BG$ have one object, endomorphism set $G$, and
composition the group multiplication. Let $i:e\to BG$ select the
object and $r:BG\to e$ be the projection. Then $ri=\id_e$.
The functor $r^*:\Set\to G\text{-}\Set$ gives the trivial action
and is fully faithful. The functor $i^*:G\text{-}\Set\to\Set$
forgets the action. On the regular two-element $G$-set $R$, there are
two equivariant endomorphisms but four underlying functions, so $i^*$
is not full. Thus $r,ri\in\Wone$ while $i\notin\Wone$.

For cohomology, the regular action has $R^G=\varnothing$, whereas
its pullback along $i$ has two global sections. For homology,
$R/G$ is a singleton whereas the pulled-back colimit is the two-element
set. Projection $r$ preserves the limits and colimits of every
coefficient on $e$, since $BG$ is nonempty and connected. Thus the same
two-out-of-three counterexample applies to both classes. The program
\texttt{code/A5\_finite\_counterexample.py} checks the category law,
all four functions, and both comparisons; the proof of full
faithfulness of trivial actions is independent of a finite size bound.
\end{proof}

\begin{claim}{A.6}{Known}{Locally constant coefficient derivators and shifts}
For every derivator $\D$ and small $A$, the prescription
\[
 \mathbb L_A(I)=
 \{F\in\D(A\times I):F\text{ is locally constant in the }A
 \text{ direction}\}
 \tag{A.6.1}
\]
defines a derivator. Its inclusion in $\D^{A}$ has both adjoints as
morphisms of derivators. Restriction along $f:A\to B$ gives a
morphism $\mathbb L_B\to\mathbb L_A$ with both adjoints.
It is fully faithful, or an equivalence, at every shape $I$ if and
only if it is so at $e$.
\end{claim}
\begin{proof}
The first two assertions are
\citet[Theorem~3.13]{Cisinski2009LocallyConstant}, translated from his covariant
notation. Compose the adjoints to the inclusions with the shifted Kan
extensions along $f\times\id_I$ to obtain the two adjoints to
restriction. These are morphisms of derivators, so their units and
counits commute with evaluation at $i\in I$. Full faithfulness of a
right adjoint is equivalent to its counit being invertible; full
faithfulness of a left adjoint is equivalent to its unit being
invertible. At every $i$, the relevant unit or counit is the one at
$e$. Der~2 therefore proves invertibility at $I$. Using both unit and
counit proves the equivalence assertion. No inference from an
arbitrary equivalence of ordinary categories to an enhancement is
being made: the specified adjunctions of derivators provide the
comparison maps.
\end{proof}

\begin{claim}{A.7}{Proved}{Finite products, coproducts, opposites, and a limit}
For a derivator, all five classes are invariant under natural
transformations between the indexing functors and under passage to
opposite functors. They are closed under finite products of arrows.
The classes $\Wlc$, $\Wone$, $\Wcoh$, and $\Whom$ satisfy the disjoint-union equivalence
\[
 \coprod_{j\in J}f_j\in\W
 \quad\Longleftrightarrow\quad (\forall j)\ f_j\in\W
 \tag{A.7.1}
\]
for every small family. For $\Wc$ only the right-to-left implication
holds in this generality; its converse is false. There is no general closure under infinite
products of arrows, even for $\Winfty$ or the set-valued classes.
\end{claim}
\begin{proof}
A natural transformation evaluated on a locally constant coefficient
is pointwise invertible; Der~2 makes it invertible in the diagram
category. The same applies to constants. For opposites, use the
twisted-arrow category $\operatorname{Tw}(A)$ with objects $a\to b$
and arrows given by pairs $a'\to a$, $b\to b'$ whose composite is
$a'\to b'$. Its two endpoint projections to $A^{\op}$ and $A$ induce
weak homotopy equivalences: its nerve is the edgewise subdivision of
$NA$, and either endpoint map is homotopic to the subdivision
homeomorphism after realization. The natural squares for $f$ compare
$f$ and $f^{\op}$ by restriction equivalences, by A.2. Such squares
preserve full faithfulness, equivalence, and the constant Hom
comparisons. For (co)homology they preserve the specified maps by
the adjunction characterization in A.2.

For $\Wlc$ and $\Wone$, A.6 gives the result for $f\times\id_I$:
inside the shift, impose local constancy also in $I$. For essential
surjectivity the preimage is locally constant in $I$ because $f^*$
reflects the invertibility of its $I$-arrows. For $\Wc$, apply the
constant homotopy-colimit characterization from A.2 and Fubini:
iterated homotopy colimits of constants over $A\times I$ identify with
the homotopy colimit over $A$ of the coefficient $\pi_{I!}\pi_I^*M$.
For $\Whom$, first integrate a locally constant coefficient over
$I$; the result is locally constant on $B$, since Kan extension in
one factor is computed pointwise in the other. Fubini then reduces
the assertion to $f$. The dual argument treats $\Wcoh$. Factoring a
finite product of arrows into these products with identities proves
the claim.

Der~1 identifies the relevant categories over a disjoint union with
products. For $\Wlc$, an equivalence of such product restriction
functors reflects each factor: fill the other factors with terminal
objects and use their unique endomorphisms for full faithfulness;
essential surjectivity is tested by filling the other target factors
with their terminal objects. This also proves the full-faithfulness
case. For cohomology choose the terminal coefficient on every other
component, and for homology choose the initial coefficient; limits,
respectively colimits, then isolate the desired comparison. For
$\Wc$, products of bijections prove the forward closure under sums,
but bijectivity of a product can conceal an empty factor. The
represented derivator of the complete lattice $\{0<1\}$ gives the
explicit failure: $\varnothing\to e$ is not in $\Wc$, whereas its
sum with $\id_e$ is in $\Wc$, since both its source and target are
nonempty. The calculation is given in G.2.

For the infinite-product failure, let $I_n$ be the path poset on
$\{0,\ldots,n\}$ with alternating orientations of consecutive
edges. There are no composable nonidentity arrows, and $BI_n$ is an
interval. Each $I_n\to e$ is a weak homotopy equivalence. In
$\prod_{n\geq1}I_n$, the objects $(0)_n$ and $(n)_n$ cannot be joined
by a finite zigzag: a zigzag of length $m$ changes the $n$th coordinate
by at most $m$, contradicting $n>m$. Its projection to $e$ fails even
bijectivity on components. This explicitly checks the infinite-family
edge case rather than transferring an untrue product rule from spaces
to nerves of categories.
\end{proof}

\begin{claim}{A.8}{Proved}{What the ordinary-category comparison does prove with an enhancement}
If $\mathcal C$ has all small limits and colimits, then
\[
 \LC_{\D_{\mathcal C}}(A)\simeq
 \Ho\Fun(BA,\mathcal C).
 \tag{A.8.1}
\]
For the specified restriction map, equivalence or full faithfulness
on these homotopy categories is equivalent to the corresponding
infinity-categorical property.
\end{claim}
\begin{proof}
The localization $NA\to\Pi_\infty A$ inverting every edge is
universal for functors carrying those edges to equivalences. Its
universal property is fully faithful on functor infinity-categories,
with exactly this image. Passing to homotopy categories preserves
that full image. The groupoid $\Pi_\infty A$ is the space $BA$;
inversion of paths implements the variance convention.

Restriction has both Kan adjoints. Its homotopy-category adjunctions
are obtained from these by applying $\Ho$. A functor is fully
faithful precisely when the appropriate adjunction unit or counit
is an equivalence. This can be tested in the homotopy category.
Equivalence is tested by both unit and counit. Equivalently, one
may recover all mapping spaces from maps to cotensors by arbitrary
small spaces, as in C.1. This argument does not say that a bare
ordinary category determines its infinity-category enhancement.
\end{proof}
